Karl Khouri's three FastAPI processes run on ports 8007--8009 and expose the required routes under \texttt{/api/Karl-Khouri}. They reuse the team's bearer-token verification and the shared SQLite database, so audit and playbook records hold foreign keys to existing users. The source files are \texttt{karl/audit\_logs.py}, \texttt{karl/playbooks.py}, \texttt{karl/reports.py}, \texttt{karl/common.py} and \texttt{karl/schema.sql}. The request and response contract is documented in \texttt{docs/karl-api.md}.

\subsubsection{Security Audit Logging (Karl Khouri)}

An audit entry records an event ID, actor, action, description, severity, related service, subject ID, the time the event occurred and the time it was stored. Severity is Information, Warning, Error or Critical. A \texttt{source} field records whether the entry came from a peer service, an administrator or the audit service itself.

Peer services send events with an \texttt{X-Service-Key} header holding \texttt{AUDIT\_SERVICE\_AUTH\_SECRET}. This key must differ from the notification key, so a credential meant for notifications cannot write audit records. Service events must carry an event ID and related service. Repeating an event ID returns the entry already stored instead of creating a second one, which makes delivery retries safe. This holds even when two deliveries of the same event arrive at once. An administrator can also add a manual entry, but the actor and time are taken from the token and the server clock. A request that tries to name a different actor or set its own time is rejected with 422.

Events are produced automatically by the other services: successful and failed logins, user creation, update and deletion, role and status changes, asset and notification changes, incident, threat and vulnerability changes, playbook activity, report generation and every request refused for lack of permission (403). Peer services first store each event in the transactional \texttt{outbox} table, then deliver it over HTTP. If the audit service is unavailable, the event stays queued for retry and the original request still succeeds. A failed login has no authenticated actor, so it is stored with a null actor ID and the attempted username as the label.

Administrators and auditors can list every entry. Administrators, analysts and auditors can retrieve a single entry and search, as the handout allows analysts to retrieve and search audit logs. Search combines actor, action, severity, related service, subject and a start/end date range. Any user can read their own activity through 	exttt{GET /logs/user/\{user\_id\}}; reading another user's activity requires the administrator or auditor role. Only administrators can delete an entry, and each deletion writes a permanent \texttt{audit\_log\_deleted} entry naming the removed record. These deletion records cannot themselves be deleted (409).

Entries cannot be edited. A database trigger rejects any UPDATE of an audit row. The only change it allows is the foreign key setting \texttt{actor\_user\_id} to null when the user is deleted. \texttt{original\_actor\_user\_id} and \texttt{actor\_label} keep the history readable after that. The audit service writes its own denials and deletions directly to its table instead of sending them to itself over HTTP, so logging cannot loop.

\subsubsection{Security Playbook and Response Automation (Karl Khouri)}

A playbook has a unique name, description, category, severity, creator and timestamps. Categories are Malware, Phishing, Ransomware, Unauthorized Access, Data Breach, Denial of Service, Insider Threat, Vulnerability Remediation and Other; severity is Low, Medium, High or Critical. The creator is taken from the token. Only administrators create and delete playbooks, as the handout lists creation for administrators only. Administrators and analysts add response actions, and analysts and auditors can read playbooks.

Each response action is an ordered step with a title, description and assignee. Steps must be numbered 1, 2, 3 and so on without gaps: a gap returns 422, and reusing a number returns 409. The database also enforces one step per number in each playbook. The assignee must be an active Security Analyst or Security Administrator. Creating a step sends that person a notification through the outbox, and retries cannot duplicate it.

A step's status moves only from Pending to In Progress to Completed. Skipping In Progress or moving backwards returns 409, and a step cannot start until every earlier step is Completed. Only the assigned user or an administrator can change a step. The start time, completion time and completing user are stored on the step. Each change also adds a row to \texttt{action\_history} with the old and new status, the user and an optional note. A playbook's detail shows its steps in order and an overall execution state of Not Started, In Progress or Completed. A playbook with a step in progress cannot be deleted. After deletion, its execution history remains readable through \texttt{GET /playbooks/\{id\}/history}.

\subsubsection{Security Reports (Karl Khouri)}

Six reports are available to administrators and auditors: summary, incidents, assets, threats, vulnerabilities and activity. Reports do not read other services' tables. Each one calls the owning service's list API over HTTP and forwards the requesting user's own token, so the upstream service applies its own permissions again. Upstream addresses are configured by environment variables, not hard-coded table access.

The incident report gives totals by status and severity, the number still open and the number unassigned. The asset report groups by type, criticality, status and department and counts active High or Critical assets. The threat report groups by type and severity. The vulnerability report groups by severity and status and gives the open count and average CVSS score. The activity report, built from the audit log, counts events by action, severity and service, failed logins, denied requests and the most active users. Every report accepts optional \texttt{start\_date} and \texttt{end\_date} filters and records a \texttt{report\_generated} audit event.

If an upstream service is unreachable, the report returns 503; if it replies with an error or an unexpected format, the report returns 502. No zero values are substituted. The summary report lists each source it could not reach under \texttt{unavailable}, shows that section as null, and sets \texttt{complete} to false. A complete summary therefore shows that every figure came from real stored records.

\subsubsection{Recorded checks and captured results (Karl Khouri)}

The saved focused test output records the original seven passing tests for the three services (\path{docs/evidence/karl-pytest-2026-10-18.txt}); two regression tests were added on 20 October 2026. They cover audit delivery from the outbox, deduplication, forged keys, actor spoofing, immutability, user deletion, playbook ordering and transitions, and report reconciliation and upstream failure. A saved six-process live HTTP check (\path{docs/evidence/karl-live-http-2026-10-18.txt}) and a saved Newman run of Karl's exported collection against nine local service processes (\path{docs/evidence/karl-postman-newman-2026-10-20.txt}) record the services working over real HTTP. The Newman run had 50 requests, 67 assertions and zero failures. Capturing the screenshots exposed one defect: after a playbook was deleted, SQLite reused its step IDs, the new step's notification event ID collided with the old one, and creating the next playbook failed with HTTP 500. AUTOINCREMENT keys and a startup migration fixed it, and a regression test covers it. The Postman and database screenshots appear in the evidence appendix.
